\documentstyle[%


righttag%


,11pt%


,amssymb%


,russian%               remove in Eng.


,izvru%                 Set izven% in Eng.


% ,izvru-ks             for brief communications


]{amsart}





% 


\newcommand{\graph}{\operatorname{graph}}


\newcommand{\co}{\operatorname{co}}


\newcommand{\const}{\operatorname{const}}


\newcommand{\meas}{\operatorname{meas}}


\newcommand{\var}{\operatornamewithlimits{var}}


\newcommand{\ess}{\operatornamewithlimits{ess}}


\newcommand{\tts}[1]{}





%\hoffset=-15mm%                  For Eng.


\setcounter{page}{19}


\god{2001}


\nomer{2~(465)} %                        Remove in Eng.


%\nomer{Vol.\,45, No.\,2}%               Set in Eng.


%\str{\pageref{begin}--\pageref{end}}%  Set in Eng.


\udk{517.977}





\author{\it В.А.\,Дыхта, О.Н.\,Самсонюк}


% NB:   ^^ remove in Eng.


\title{Принцип максимума в негладких задачах оптимального импульсного


управления с многоточечными фазоограничениями}


\thanks{Работа выполнена при финансовой поддержке Российского фонда


фундаментальных исследований (грант


\No\,98-01-00837).}





\date{}


%\pagestyle{myheadings}%         Set in Eng.


\pagestyle{plain} %              Remove in Eng.


\begin{document}


%\thispagestyle{plain}%          Set in Eng.


\maketitle


\label{begin}





\section{Введение}





В работе получены необходимые условия


оптимальности первого порядка в форме принципа максимума (ПМ) в


задаче оптимального управления с траекториями ограниченной вариации


при наличии общих многоточечных фазоограничений.


Задача оптимального импульсного управления рассматривается


без предположений о гладкости на входные данные ---


предполагается измеримая зависимость по времени и липшицевая по


фазовым координатам. Кроме того, не предполагается выполнение


условия корректности, что приводит к неоднозначности реакции


динамической системы на импульсное управление (векторную меру); как


следствие, импульсному управлению и начальному условию может


соответствовать воронка обобщенных решений.





Приведем определение обобщенного (об.) решения, устоявшееся в теории


оптимальных импульсных процессов \cite{80}, \cite{75}.


Рассмотрим управляемую систему


\begin{gather}


\Dot{x}=f(t,x,V,u)+G(t,x,V)v, \quad


\Dot{V}=\|v\|, \label{eq:(5.1)} \\


u(t) \in U,\quad v(t) \in K,\quad t \in [t_0,t_1].


\label{eq:(5.2)}


\end{gather}


Здесь $x(\cdot)$, $V(\cdot)$ ---


абсолютно непрерывные, $u(\cdot)$, $v(\cdot)$ --- измеримые


ограниченные функции, $U$ --- компактное множество в $R^{d(u)}$, $K$


--- выпуклый замкнутый конус в $R^{d(v)}$,


$\|v\|=\sum\limits_{j=1}^{d(v)}|v_j|$,


$d(z)$ обозначает размерность вектора $z$.


Если не оговорено противное, то термины ``измеримость''


и ``ограниченность'' применительно к


функциям относятся к мере Лебега ${\cal L}$, а все соотношения,


содержащие измеримые функции, считаются выполненными ${\cal L}$-почти всюду.


\medskip





{\bf Определение} (\cite{80}). Пара


функций $x(\cdot)$, $V(\cdot)$, непрерывных справа на $(t_0,t_1]$ и


имеющих ограниченную вариацию, называется обобщeнным


решением системы (\ref{eq:(5.1)}), (\ref{eq:(5.2)}), если существует


такая последовательность функций $\{x_n(\cdot)$, $V_n(\cdot),$


$u_n(\cdot)$, $v_n(\cdot)\}$, удовлетворяющих системе


(\ref{eq:(5.1)}), (\ref{eq:(5.2)}), что


$$


\sup_n\|v_n(\cdot)\|_{L_1} < \infty, \quad (x_n,V_n) \rightarrow


(x,V) \ \; \text{слабо}^{\ast}\ \text{в}\ \; BV


$$


(тогда


$(x_n(t),V_n(t)) \rightarrow (x(t),V(t))$ в точках непрерывности


$(x,V)$ и в концах отрезка $T$).


\medskip





Если положить


$w_n(t)=\int\limits_{t_0}^{t}v_n(\tau)d\tau,$


то из теоремы Хелли


следует, что последовательность $\{ w_n(\cdot)\}$ можно


считать слабо$^{\ast}$ сходящейся к некоторой функции


ограниченной вариации $w(\cdot)$. При этом порожденная ею мера


Лебега--Стильтьеса $d w$ оказывается $K$-значной вследствие


выпуклости и замкнутости конуса $K$. Заметим также, что хотя


$V_n(t)=\var\limits_{[t_0,t]} w_n(\cdot),$


предельная функция


$V(\cdot)$ не обязана совпадать с вариацией вектор-функции


$w(\cdot)$.





В \cite{80} показано, что при стандартных


предположениях непрерывности, липшицевости и не более чем линейного


роста функций $f$, $G$ по $(x,V)$, а также выпуклости годографа, для


любого обобщенного решения $(x(\cdot), V(\cdot))$ найдутся


$K$-значная


мера Лебега--Стильтьеса $dw$,\linebreak $\cal{L}$- и $dV$-измеримая


ограниченная функция $u(\cdot)$, $\cal{L}$-измеримые


функции $\omega_s(\cdot)$, определенные для каждого момента $s \in


S_d(V)=\{s\mid [V(s)]:=V(s)-V(s-)>0\}$ на отрезке


$[0,d_s]=[0,[V(s)]]$, для которых выполняются условия


\begin{equation}


\omega_s(\tau) \in K_1:=\{v\in K\mid \|v\|\leq 1\}, \quad


\int_0^{d_s}\omega_s(\tau)d\tau=[ w (s)]; \label{B.4}\end{equation}


функции


$(x(\cdot),V(\cdot))$ удовлетворяют системе дифференциальных


уравнений с мерой


\begin{gather}


dx(t)=f(t,x(t),V(t),u(t))dt+G(t,x(t),V(t))d w_c(t)


+\sum\begin{Sb} s \in S_d(V) \\ s \leq t\end{Sb}


[x(s)]\delta(t-s),\label{eq:(5.4)} \\


V(t)=\var_{[t_0,t]}


w_c(\cdot)+ \sum\begin{Sb} s\in S_d(V) \\ s \leq t\end{Sb}[V(s)],\quad


V(t_0-)=0, \label{eq:(5.5)}


\end{gather}


где $ w_c$ --- непрерывная составляющая в


разложении Лебега функции $ w$; скачки об. решения в точках $s\in


S_d(V)$ находятся из условий


\begin{equation}


[x(s)]=z_s(d_s)-x(s-), \quad {[V(s)]} = z_{Vs}(d_s)-V(s-),


\label{eq:(5.6)}


\end{equation}


где функции $z_s(\tau)$, $z_{Vs}(\tau)$


являются решениями предельной системы


\begin{equation}


\begin{array}{c}


\displaystyle\frac{dz_s}{d\tau}=G(s,z_s,z_{Vs})\omega_s(\tau),


\quad z_s(0)=x(s-),\\[3mm]


\displaystyle\frac{dz_{Vs}}{d\tau}=1, \quad z_{Vs}(0)=V(s-), \quad


\tau\in [0,d_s].


\end{array}


\label{eq:(5.7)}


\end{equation}





Таким образом, индивидуальная об. траектория является ветвью


интегральной воронки об. решений, соответствующей $(x_0,u(\cdot),dw


)$. Эта ветвь определяется вспомогательными предельными управлениями


$\omega_s(\tau)$, которые отражают способ приближения


последовательности $\{ w_n\}$ к функции $ w$ и описывают скачок


траектории в моменты разрыва функции $V$. Обобщенный, импульсный


процесс системы (\ref{eq:(5.1)}), (\ref{eq:(5.2)})


представляет собой мультинабор


$$


e=[x(\cdot), V(\cdot),u(\cdot), dw,


\{z^i(\cdot),z_V^i(\cdot),\omega^i(\cdot)\}_{s_i\in S_d(V)}],


$$


в котором $z^i(\cdot)=z_{s_i}$, $s_i\in S_d(V)$, и


аналогично определены функции $z_V^i$, $\omega^i$.





В данной работе мы используем приведенное описание множества об.


решений при ослабленных предположениях относительно функций $f$ и $G$,


а множество $U$ считаем переменным. Это естественно, 


поскольку при доказательстве необходимых условий оптимальности факт


существования исследуемого решения постулируется, так что условия типа


непрерывности по $t$ и роста по $(x,V)$ становятся излишними. Кроме


того, ослабление требований касается в основном функции $f$, свойства


которой при переходе к импульсному линейному управлению играют


второстепенную роль и в действительности допускают ослабление


(\cite{75}, сc.\,173, 182). Заметим также, что


доказательство представленного ПМ не использует условий аппроксимируемости.





Отметим, что наиболее общий ПМ для гладких задач оптимального


импульсного управления при невыполнении условия корректности получен


в \cite{80} методом разрывной замены времени. Необходимые условия


оптимальности при негладкой зависимости входных данных получены в


\cite{83}, \cite{Diz.vuz} для более частных вариантов задачи (при


отсутствии многоточечных фазоограничений и выполнении условия


корректности); в основу их доказательства положен вариационный


принцип Экланда. В \cite{DS99} доказан вариант ПМ для


негладкой задачи оптимального управления с траекториями ограниченной


вариации на фиксированном отрезке времени без промежуточных


фазоограничений, техника его доказательства распространяется на


более сложный класс задач данной статьи.





\section{Постановка задачи и предположения}








Пусть $\theta=(\theta_0,\theta_1,\dotsc,\theta_n)$ ---


вектор нефиксированных моментов времени таких, что $\theta_0\leq


\theta_1\leq \dotsb \leq \theta_n$, $n< \infty$,


$x(\theta)$, $V(\theta)$, $x(\theta-)$, $V(\theta-)$ --- векторы,


составленные из односторонних пределов в моменты $\theta_j$


непрерывных справа функций ограниченной вариации $x(\cdot)$,


$V(\cdot)$, т.\,е.


\begin{alignat*}{2}


x(\theta)&=(x(\theta_0),\dotsc,x(\theta_n)),&\quad


x(\theta-)&=(x(\theta_0-),\dotsc,x(\theta_n-)),\\


V(\theta)&=(V(\theta_0),\dotsc,V(\theta_n)),


&\quad V(\theta-)&=(V(\theta_0-),\dotsc,V(\theta_n-)).


\end{alignat*}





Рассматривается задача ${\cal P}(\theta)$ минимизации функционала


$$


l(b), \quad b=(\theta,x(\theta-),x(\theta),V(\theta-),V(\theta)),


$$


на множестве импульсных процессов


$e=[\theta, x(\cdot), V(\cdot), u(\cdot), d w,


\{z^i(\cdot),z_V^i(\cdot),\omega^i(\cdot)\}_{s_i\in S_d(V)}],$


удовлетворяющих следующим условиям допустимости:


\begin{itemize}


\item[1)] $x(\cdot)$, $V(\cdot)$ --- непрерывные справа на


$(\theta_0,\theta_n]$ функции ограниченной вариации;


\item[2)] $u(\cdot)$ --- измеримая функция, удовлетворяющая ограничению


$$


u(t) \in U(t),\quad t\in [\theta_0,\theta_n],


$$


где $U(t)$ --- сечение множества $U{\subset} R{\times} R^{d(u)}$ при


фиксированном $t$, т.\,е.


$U(t){=}\{u\mid (t,u){\in} U\};$


\item[3)] $d w(t)$ --- $K$-значная мера Лебега--Стильтьеса на


$[\theta_0,\theta_n]$, т.\,е. $ d w(B) \in K$ для любого борелевского


множества $B \subset [\theta_0,\theta_n]$;


\item[4)] компоненты набора $e$ удовлетворяют ограничениям


(\ref{B.4})--(\ref{eq:(5.7)});


\item[5)] выполняется многоточечное фазоограничение


\begin{equation}


b\in C. \label{eq:(5.44)}


\end{equation}


\end{itemize}





Через $\ov e=


[\ov \theta, \ov


x(\cdot),\ov V(\cdot),\ov u(\cdot),d\ov w,


\{\ov z^i(\cdot),\ov z_V^i(\cdot),\ov\omega^i(\cdot)\}_{s_i\in


S_d(\ov V)}]$ будем обозначать


допустимый импульсный процесс, исследуемый на оптимальность.





Поставленная задача


рассматривается при следующих предположениях:


\begin{itemize}


\item[(B1)] множество $S_d(\ov{V})$ конечно и мера


$d\ov{ w}$ не имеет непрерывной сингулярной составляющей, т.\,е.


$$


d\ov{ w}(t)=\ov{v}(t)dt+\sum_{\ov s_i \in


{S}_d(\ov V)}\ov{c}^i\delta(t-\ov s_i),


$$


причем $\ov v\in


L_\infty ([\ov\theta_0,\ov\theta_n],R^{d(w)});$


\item[(B2)] множество $U\subset R\times R^{d(u)}$ борелевское, $K


\subset R^{d(v)}$ --- замкнутый выпуклый конус, $C$ --- замкнутое


множество;


\item[(B3)] при всех $(x,V)$ функция $(t,u) \rightarrow f(t,x,V,u)$


\; ${\cal L}\times {\cal B}$-измерима;


\item[(B4)] существуют окрестность ${\cal Q}$ множества


$$


\ov\graph\{\ov{x}(\cdot),\ov{V}(\cdot)\}


=\{(t,x,V)\mid


t \in \Delta, \ x \in [\ov{x}(t-),\ov{x}(t+)],\ V \in


[\ov{V}(t-),\ov{V}(t+)]\},


$$


где $\Delta\supset [\ov\theta_0,\ov\theta_n]$


и число $M>0$ такие, что


выполняются условие Липшица


$$


|f(t,x,V,u)-f(t,x',V',u)| \leq


M(|x-x'|+|V-V'|)\ \; \forall\, (t,x,V,u), (t,x',V',u) \in {\cal


Q}\times U(t)


$$


и условие ограниченности


$$


|f(t,x,V,u)| \leq M \ \;


\forall\, (t,x,V,u) \in {\cal Q}\times U(t);


$$


\item[(B5)] матрица $G$ непрерывно дифференцируема;


\item[(B6)] функция $l$ локально липшицева;


\item[(B7)] если $\ov\theta_{i}\in S_d(\ov


V)$, то $\ov\theta_i\ne\ov\theta_{j}$ при $j\ne i$.


\end{itemize}





{\bf Замечания.} 1) Предположение (B1) необходимо для


используемого метода доказательства, оно позволяет охватить большинство


приложений.





2) Предположение (B3) выполняется, если, например, функция $f$


измерима по $t$ и непрерывна по $u$.





3) Множество $\ov\graph\{\ov{x}(\cdot),\ov{V}(\cdot)\}$ получается


из графика разрывной вектор-функции


$(\ov{x}(\cdot),$ $\ov{V}(\cdot))$ путем его дополнения


вертикальными отрезками, соединяющими односторонние пределы в точках


разрыва.





4) В задаче ${\cal P}(\theta)$ односторонние пределы об. траекторий в


нефиксированные моменты $\theta_j$ стеснены общим ограничением


(\ref{eq:(5.44)}) (для симметрии от условия $V(\theta_0-)=0$ мы


отказываемся). Моменты $\theta_j$ могут совпадать с действием


импульсов, но предположение (B7) исключает в этом случае совпадение


нескольких компонент вектора $\theta$ с моментом импульса.


\medskip





Таким образом, ${\cal P}(\theta)$ --- это некоторая задача оптимизации


с разрывными траекториями при наличии фазоограничений в промежуточные


и конечные моменты времени. Конечно, выполнение этих ограничений на


об. решении не обеспечивает выдерживание их на аппроксимирующих


обычных решениях: таких аппроксимаций может вовсе не существовать уже


в случае более простых терминальных ограничений


\cite{Diz.vuz}, \cite{DIFAC} или же соблюдение ограничения


(\ref{eq:(5.44)}) оказывается возможным лишь на некоторых


аппроксимациях. Рассматриваемый ниже пример 1 иллюстрирует данное


свойство. Он показывает также, что проекция ограничения


(\ref{eq:(5.44)}) в чисто фазовое ограничение для предельных


подпроцессов приводит к совершенно другой задаче оптимального


импульсного управления.





\section{Принцип максимума}





Введем необходимые обозначения и понятия. Пусть


\begin{gather*}


H(t,x,V,\psi,u,v )=\langle \psi,f(t,x,V,u)+G(t,x,V)v\rangle ,\\


H_{0}(t,x,V,\psi,u )=\langle \psi,f(t,x,V,u)\rangle ,\quad


H_1(t,x,V,\psi,v)=\langle \psi,G(t,x,V)v\rangle, \\


{\cal H}_0(t,s)= \sup_{u \in U(t)}


H_0(t,\ov{x}(s),\ov{ V}(s),\psi(s),u).


\end{gather*}





Если $\varphi (t,x,V,\psi,v)$ --- заданная функция, то через


$\varphi|_{\Theta(s_i)}(s_i,z^i,z_V^i,p^i,\omega^i)$


будет обозначаться замена переменных $(t,x, V,\psi,v)


\rightarrow (s_i,z^i,z_V^i,p^i,\omega^i)$ в функции


$\varphi$. Запись типа $H_0(t)$ будет означать вычисление функции


$H_0$ вдоль исследуемого процесса $\ov{e}$, а запись $H_0(t,u)$ ---


фиксирование всех опущенных аргументов функции $H_0$ вдоль процесса


$\ov{e}$ (в данном случае $x=\ov{x}(t)$, $V=\ov


V(t)$, $\psi=\psi(t)$, где $\psi(t)$ --- траектория двойственных


переменных). Точно так же следует понимать обозначение


$\varphi|_{\Theta(s_i)}(\tau)$.





Поскольку моменты действия импульсов для


допустимых процессов не фиксируются, а зависимость функции $f$ по $t$


лишь измеримая, то для корректной записи условий оптимальности


моментов скачка траектории используем следующее понятие из теории


функций \cite{Clarke.Vinter}.





Пусть $I$ --- интервал из $R$, $s \in I$ --- заданная точка и $\varphi: I


\rightarrow R^n$ --- измеримая функция. Определим множество


$\ess\limits_{t \rightarrow s}\varphi (t)$ существенных


значений функции $\varphi$ в точке $s$


как совокупность всех векторов


$\eta\in R^n$, для которых при любом $\varepsilon > 0$


$$


{\cal L}\text{-}\meas\{t\mid s-\varepsilon <t< s+\varepsilon,\


|\eta-\varphi(t)|< \varepsilon \}>0.


$$


Легко проверить, что если $s$ --- точка непрерывности $\varphi$, то


$\ess\limits_{t \rightarrow s}\varphi(t)=\{\varphi(s)\}$,


а если $s$ --- точка разрыва 1-го рода, то


$\ess\limits_{t \rightarrow s}\varphi(t)= \{\varphi(s-),


\varphi(s+)\}$.





Если $g(x,y)$ --- липшицева функция, то символ


$\partial g(\ov{x},\ov{y})$ обозначает обобщенный


градиент Кларка по $(x,y)$, вычисленный в точке $(\ov{x},\ov{y})$


(\cite{23}, с.\,17). Символ $\co C$ означает


выпуклую оболочку множества $C$, а $N_C(b)$ --- нормальный конус к $C$ в


точке $b$ (\cite{23}, c.\,54).





Сформулируем принцип максимума.





\begin{thm}


Пусть $\ov e$ --- оптимальный процесс в задаче


${\cal P}(\theta)$. Тогда существуют число $\alpha_0\in \{0,1\}$,


числа $h_s^{-}$, $h_s^{+}$,


определенные для каждого момента $s\in


\{\ov\theta_0,\dotsc,\ov\theta_n\}$,


функции


$(\psi,\psi_V):[\ov\theta_0,\ov\theta_n]\rightarrow


R^{d(x)}\times R$ и функции $(p^i,p_V^i):[0,\ov d_i]\rightarrow


R^{d(x)}\times R$, определенные для каждого момента $\ov s_i\in


S_d(\ov V)$, где $\ov d_i=[\ov V(\ov s_i)]$,


такие, что выполнены следующие условия$:$





$1)$ условие нетривиальности


$$


\alpha_0+ |\wt \psi(\ov \theta_n+)|+\sum_{j=0}^n|\wt


\psi(\ov \theta_j-)|


+\sum_{s_i\in S_d(\ov V)\cap


\{\ov\theta_0,\dotsc,\ov\theta_{n}\}}|\wt


p^i([\ov V(s_i)])|>0,


$$


где $\wt\psi=(\psi,\psi_V)$, $\wt p^i=(p^i,p_V^i)$;





$2)$ функции $\psi(\cdot)$, $\psi_V(\cdot)$ кусочно абсолютно


непрерывны, на интервалах абсолютной непрерывности удовлетворяют


сопряженному дифференциальному включению


\begin{equation}


\big(\Dot{\psi}(t)+H_{1x}(t),\ \Dot{\psi}_V(t)+H_{1V}(t)\big) \in


-\partial_{(x,V)} {H}_0(t), \label{eq:(5.45)}


\end{equation}


функции $p^i(\cdot)$, $p_V^i(\cdot)$ абсолютно непрерывны,


удовлетворяют дифференциальным уравнениям


\begin{equation}


\frac{dp^i(\tau)}{d\tau}= -


{H}_{1x} \big|_{\Theta(\ov{s}_i)}(\tau), \quad


\frac{dp_V^i(\tau)}{d\tau}= -


{H}_{1V} \big|_{\Theta(\ov{s}_i)}(\tau) \label{eq:(5.46)}


\end{equation}


и связаны с $\psi(\cdot)$,


$\psi_V(\cdot)$ в моменты $\ov s_i\in S_d(\ov V)\setminus


\{\ov\theta_0,\dotsc,\ov\theta_n\}$ условиями допустимости скачка


\begin{alignat*}{2}


\psi(\ov s_i+)&=p^i(\ov d_i),&\quad


\psi(\ov s_i-)&=p^i(0),\\


\psi_V(\ov s_i+)&=p_V^i(\ov d_i),&\quad \psi_V(\ov s_i-)&=


p_V^i(0);


\end{alignat*}





$3)$ условия максимума


\begin{gather}


{H}_0(t)= \max_{u \in U(t)}


H_0(t,\ov{x}(t),\ov{V}(t),\psi(t), u), \quad t


\in [\ov \theta_0,\ov \theta_n],\notag\\ 


%


\langle {H}_v(t),\ov{v}(t)\rangle


+\psi_V(t)\|\ov{v}(t)\|= 0=


\max_{v \in K}(\langle {H}_v(t),v\rangle


+\psi_V(t)\|v\|), \quad \ t \in [\ov


\theta_0,\ov \theta_n], \label{eq:(5.9)} \\


\langle {H}_{v} \big|_{


\Theta(\ov{s}_i)}(\tau),\ov{\omega}^i(\tau)\rangle


+p_V^i(\tau)= 0= 


\max_{\omega \in K_1}\langle {H}_{v}


\big|_{ \Theta(\ov{s}_i)}(\tau),\omega\rangle + p_V^i(\tau),


\quad \tau \in [0,\ov d_i],\ \;\overline s_i\in S_d(\ov


V); \label{eq:(5.10)}


\end{gather}





$4)$ условия оптимальности моментов $\theta_j$ и $s_i:$





\begin{gather*}


h_s^{+}\in


\co \ess_{t\rightarrow s}{\cal H}_0(t, s+),\quad


h_s^{-}\in


\co \ess_{t\rightarrow s}{\cal H}_0(t, s-)


\ \; \forall s\in \{\ov\theta_0,\dotsc,\ov\theta_n\}, \\


%


\int_0^{[\ov V(s)]} {H}_{1t} \big|_{ \Theta(s)}(\tau)d\tau


\in \co \ess_{t\rightarrow {s}} {\cal


H}_0(t,{s}+) - \co\ess_{t\rightarrow


{s}} {\cal H}_0(t,{s}-)\ \;


\forall s\in S_d(\ov V)\setminus


\{\ov\theta_0,\dotsc,\ov\theta_n\};


\end{gather*}





$5)$ условие трансверсальности


$$


\big(\eta_\theta,\eta_{x(\theta-)},\eta_{x(\theta)},


\eta_{V(\theta-)},\eta_{V(\theta)}\big)\in \alpha_0\partial


l(\ov b)+N_C(\ov b),


$$


где $\ov b=(\ov\theta,\ov x(\ov\theta-),\ov


x(\ov\theta),\ov V(\ov\theta-),\ov V(\ov\theta)),$


\begin{align*}


\eta_{\theta}&=((-h_{\ov


\theta_0}^{+}+r(\ov\theta_0)),\


(h_{\ov\theta_j}^{-}-h_{\ov\theta_j}^{+}+r(\ov


\theta_j))_{j=\ov{1,n-1}},


(h_{\ov\theta_n}^{-}+r(\ov


\theta_n))),\\


%


\eta_{x(\theta-)}&=(\psi(\ov\theta_0-),\


(-\psi(\ov \theta_j-)+


p_{\ov\theta_j}(0))_{j=\ov{1,n}}), \\


%


\eta_{x(\theta)}&=((\psi(\ov \theta_j+)-


p_{\ov\theta_j}( [\ov


V(\theta_j)]))_{j=\ov{0,n-1}}, -\psi(\ov\theta_n+)), \\


%


\eta_{V(\theta-)}&=(\psi_V(\ov\theta_0-), (-\psi_V(\ov \theta_j-)+


p_{V\ov\theta_j}(0))_{j=\ov{1,n}}), \\


%


\displaybreak[3]


\eta_{V(\theta)}&=((\psi_V(\ov \theta_j+)-


p_{V\ov\theta_j}( [\ov


V(\theta_j)]))_{j=\ov{0,n-1}}, -\psi_V(\ov\theta_n+)), \\


%


\psi(\ov\theta_0-)&=p_{\ov\theta_0}(0), \quad


\psi(\ov\theta_n+)=p_{\ov\theta_n}([\ov V(\ov


\theta_n)]), \\


%


\psi_V(\ov\theta_0-)&=p_{V\ov\theta_0}(0), \quad


\psi_V(\ov\theta_n+)=p_{V\ov\theta_n}([\ov


V(\ov \theta_n)]),\\


%


p_{\ov\theta_j}(\tau)&=


\begin{cases}


p^i(\tau), & \text{если \;}


\ov\theta_j=\ov s_i\in S_d(\ov V); \\


\const, & \text{если \;} \ov\theta_j\notin S_d(\ov V),


\end{cases}


\\


%


p_{V\ov\theta_j}(\tau)&=


\begin{cases}


p_V^i(\tau), & \text{если \;}


\ov\theta_j=\ov s_i\in S_d(\ov V);\\


\const, & \text{если \;} \ov\theta_j\notin S_d(\ov V),


\end{cases}


\\


%


r(\ov\theta_j)&=


\begin{cases}


\displaystyle\int_0^{\ov d_i}H_{1t}\big|_{ \Theta(\ov


s_i)}(\tau)d\tau, & \ov\theta_j=\ov s_i\in


S_d(\ov V);\\


0, & \ov\theta_j\notin


S_d(\ov V).


\end{cases}


\end{align*}


\end{thm}





{\bf Замечания.} 1) Если в задаче ${\cal P}(\theta)$ функционал $l$ и


множество $C$, задающее многоточечные фазоограничения, не зависят от


значений $x(\theta_0)$, $V(\theta_0)$, $x(\theta_n-)$, $V(\theta_n-)$,


то дополнительно к условиям теоремы 1 выполняются условия


$$


h_{\ov\theta_0}^{-}-h_{\ov\theta_0}^{+}+r(\ov\theta_0)=-a_0,\quad


h_{\ov\theta_n}^{-}-h_{\ov\theta_n}^{+}+r(\ov\theta_n)=a_1


$$


при некоторых $a_0$, $a_1\geq 0$.





2) Если процесс $\ov e$ имеет совпадающие моменты


$$


\ov\theta_j=\ov\theta_{j+1}=\cdots=\ov\theta_{j+k},


$$


то


$h_{\ov\theta_i}^{+}=h_{\ov\theta_{i+1}}^{-}$,


$i=\ov{j,j+k-1}.$








\section{Обсуждение принципа максимума и частные случаи}





Условия трансверсальности и оптимальности промежуточных моментов


теоремы 1 выглядят громоздко, что связано с наличием общего


ограничения типа включения на промежуточные моменты времени и значения


траектории. Ниже мы приведем более компактную формулировку ПМ для


задачи, в которой промежуточные и терминальные


фазоограничения заданы функционально.


\medskip





\begin{center}{\bf 4.1. Случай функционального задания многоточечных


фазоограничений}


\end{center}


\medskip








Рассмотрим задачу ${\cal P}_1(\theta)$, в которой множество $C$ задано


функционально в виде


$$


\varphi(b)\leq 0,\quad k(b)=0,\quad


V(\theta_0-)=0,


$$


где вектор-функции $\varphi$, $k$ имеют размерности


$d(\varphi)$, $d(k)$ соответственно и удовлетворяют


предположению


\begin{itemize}


\item[(B8)] функции $l$, $\varphi_j$, $j=\ov{1,d(\varphi)}$,


локально липшицевы и регулярны (\cite{23}, с.\,44); вектор-функция $k$


непрерывно дифференцируема.


\end{itemize}





Функциональное задание множества $C$ и локальная липшицевость функций


$l$, $\varphi$, $k$ позволяет сформулировать ПМ в терминах множителей


Лагранжа. Введем функцию Лагранжа


$$


L(b)=\alpha_0 l(b)+\langle \alpha,\varphi(b)\rangle +\langle


\beta,k(b)\rangle.


$$


Предположение (B8) обеспечивает регулярность


функции Лагранжа, что позволяет представить ПМ в более компактной


формулировке, т.\,к. в этом случае обобщенный градиент $\partial


L(b)$ содержится в прямом произведении частных обобщенных градиентов.





ПМ для задачи ${\cal P}_1(\theta)$ дает





\begin{thm}


Пусть $\ov e$ --- оптимальный процесс в задаче


${\cal P}_1(\theta)$. Тогда существуют число $\alpha_0$,


векторы $\alpha\in R^{d(\varphi)}$, $\beta\in R^{d(k)}$, функции


$(\psi,\psi_V):[\ov\theta_0,\ov\theta_n]\rightarrow


R^{d(x)}\times R$ и функции $(p^i,p_V^i):[0,\ov d_i]\rightarrow


R^{d(x)}\times R$, $\ov d_i=[\ov V(\ov s_i)]$,


определенные для каждого момента $\ov s_i\in S_d(\ov V)$,


такие, что выполнены следующие условия$:$





$1)$ условие нетривиальности теоремы $1$, $\alpha_0\geq 0$,


$\alpha\geq 0$, $\langle \alpha,\varphi(\ov{b})\rangle =0;$





$2)$ функции $\psi(\cdot)$, $\psi_V(\cdot)$ кусочно абсолютно


непрерывны, на интервалах абсолютной непрерывности удовлетворяют


дифференциальному включению $(\ref{eq:(5.45)})$, функции


$p^i(\cdot)$, $p_V^i(\cdot)$ абсолютно непрерывны, удовлетворяют


дифференциальным уравнениям $(\ref{eq:(5.46)});$





$3)$ условия максимума теоремы $1;$





$4)$ условия скачка и трансверсальности:





{\em а)} для всех $s\in


\{\ov\theta_1,\dotsc,\ov\theta_{n-1}\}\setminus S_d(\ov V)$


\begin{align*}


\psi(s+)-\psi(s-)&\in


\sum_{\ov\theta_j=s}(


\partial_{x(\theta_j-)}L(\ov


b)+\partial_{x(\theta_j)}L(\ov b)), \\


\psi_V(s+)-\psi_V(s-)&\in \sum_{\ov\theta_j=s}(


\partial_{V(\theta_j-)}L(\ov


b)+\partial_{V(\theta_j)}L(\ov b));


\end{align*}





{\em б)} для всех $\ov s_i\in S_d(\ov V)\setminus\{


\ov\theta_0, \dotsc, \ov\theta_{n}\}$


\begin{alignat*}{2}


\psi(\ov


s_i+)-p^i(\ov d_i)&=0, &\quad p^i(0)-\psi(\ov s_i-)&=0, \\


\psi_V(\ov s_i+)-p_V^i(\ov d_i)&=0, &\quad


p_V^i(0)-\psi_V(\ov s_i-)&=0;


\end{alignat*}





{\em в)} для всех $\ov s_i\in S_d(\ov V)\cap\{


\ov\theta_1, \dotsc, \ov\theta_{n-1}\}$


\begin{alignat*}{2}


\psi(\ov


s_i+)-p^i(\ov d_i)&\in \partial_{x(\ov s_i)}L(\ov


b),&\quad p^i(0)-\psi(\ov s_i-)&\in \partial_{x(\ov


s_i-)}L(\ov b),\\


\psi_V(\ov s_i+)-p_V^i(\ov


d_i)&\in \partial_{V(\ov s_i)}L(\ov b), &\quad


p_V^i(0)-\psi_V(\ov s_i-)&\in \partial_{V(\ov


s_i-)}L(\ov b);


\end{alignat*}





{\em г)} \hspace{2.5cm} $\psi(\ov \theta_0+)-


% adjust in eng. ^^


p_{\ov\theta_0}( [\ov V(\theta_0)])\in L_{x(


\theta_0)}(\ov b)$,\ \ $ -\psi(\ov \theta_n-)+


p_{\ov\theta_n}( 0)\in L_{x( \theta_n-)}(\ov b)$,


\begin{align*}


-\psi_{V}(\ov \theta_n-)+


p_{V\ov\theta_n}( 0)&\in L_{V(


\theta_n-)}(\ov b), \\


(\psi(\ov \theta_0-),-\psi(\ov\theta_n+),-\psi_V(\ov\theta_n+))


&\in \partial_ {(x(\theta_0-),x(\theta_n),V(\theta_n))}


L(\ov{b});


\end{align*}





{\em 5)} условия оптимальности моментов


$\forall s\in \{\ov\theta_1,\dotsc,\ov\theta_{n-1}\}


\cup S_d(\ov V)$, $s\ne \ov\theta_0,


\ov\theta_n$,


\begin{align*}


\int_0^{[\ov V(s)]} {H}_{1t}


\big|_{\Theta(s)}(\tau)d\tau &\in \co


\ess_{t\rightarrow s} {\cal H}_0(t,{s}+) - \co


\ess_{t\rightarrow {s}} {\cal H}_0(t,{s}-) +


\sum_{\ov\theta_j=s} \partial_{\theta_j}L(\ov b), \\


%


\int_0^{[\ov V(\ov\theta_0)]} {H}_{1t}


\big|_{\Theta(\ov\theta_0)}(\tau)d\tau &\in


\co \ess_{t\rightarrow \ov\theta_0} {\cal


H}_0(t,\ov\theta_0+)


+\sum_{\ov\theta_j=\ov\theta_0} \partial_{\theta_j}L(\ov b), \\


%


\int_0^{[\ov V(\ov\theta_n)]} {H}_{1t}


\big|_{\Theta(\ov\theta_n)}(\tau)d\tau &\in -


\co\ess_{t\rightarrow \ov\theta_n} {\cal


H}_0(t,\ov\theta_n-)


+\sum_{\ov\theta_j=\ov\theta_n} \partial_{\theta_j}L(\ov


b).


\end{align*}


\end{thm}





Теорема 2 следует из теоремы 1 с учетом свойств обобщенных градиентов


и нормалей.





Выделим случай гладкой задачи ${\cal P}_{s}$, предположив


дополнительно, что функция


$f(t,x,V,u)$ непрерывно дифференцируема по переменным $t$, $x$, $V$ и


непрерывна по $u$, функции $l$, $\varphi$ и $k$


непрерывно дифференцируемые. Предположим также, что


функции $l$, $\varphi$ и $k$ не зависят от $x(\theta_0)$,


$V(\theta_0)$, $x(\theta_n-)$, $V(\theta_n-)$.





ПМ для задачи ${\cal P}_{s}$ содержит условия 1)--4) теоремы


2 с заменой включений на равенства и обобщенных градиентов на


обычные частные производные. Дополнительно в нем утверждается


существование функций


$\psi_t:[\ov\theta_0,\ov\theta_n]\rightarrow R$,


$p_t^i:[\ov\theta_0,\ov\theta_n]\rightarrow R$, $ s_i \in


S_d(\ov{V})$, где $\psi_t(\cdot)$ кусочно абсолютно


непрерывна, удовлетворяет на интервалах абсолютной непрерывности


дифференциальному уравнению


$$


\Dot{\psi}_t(t)=- {H}_t(t),


$$


а функции $p_t^i(\cdot)$ абсолютно непрерывны на $[0,\ov{d}_i]$ и


удовлетворяют соответствующим предельным сопряженным уравнениям


$$


\frac{dp_t^i(\tau)}{d\tau}=- {H}_{1t} \big|_{


\Theta(\ov{s}_i)}(\tau).


$$


При этом выполняются следующие условия скачка и трансверсальности:


\begin{gather*}


\psi_t(s+)-\psi_t(s-)=


\sum_{\ov\theta_j=s}L_{\theta_j}(\ov b)


\ \; \forall s\in\{\ov\theta_1,\dotsc,\ov\theta_{n-1}\}\setminus


S_d(\ov V), \\


%


\psi_t(\ov


s_i+)-p_t^i(\ov d_i)=0, \quad p_t^i(0)-\psi_t(\ov s_i-)=0


\ \; \forall\ \ov s_i\in S_d(\ov V)\setminus\{


\ov\theta_1, \dotsc, \ov\theta_{n-1}\}, \\


%


\psi_t(\ov s_i+)-p_t^i(\ov


d_i)=L_{\ov s_i}(\ov b),\quad p_t^i(0)-\psi_t(\ov


s_i-)=0


\ \; \forall\ \ov s_i\in S_d(\ov V)\cap\{


\ov\theta_1, \dotsc,


\ov\theta_{n-1}\}, \\


\psi_t(\ov\theta_0-)=L_{\theta_0}(\ov b),\quad


\psi_t(\ov\theta_n+)=-L_{\theta_n}(\ov b).


\end{gather*}


Условие 5)


теоремы 2 с учетом замечания 1) к теореме 1 примет вид


$$


H_0(t)+\psi_t(t)\quad


\begin{cases}


\leq 0, & {t}=\ov{\theta}_0; \\


= 0, & {t} \in (\ov{\theta}_0,\ov{\theta}_n); \\


\geq 0, & {t}=\ov{\theta}_n.


\end{cases}


$$





\begin{center}{\bf 4.2. Частные случаи условий максимума}


\end{center}


\medskip





а) Рассмотрим типичный для приложений случай $K=R_{+}^{d(v)}$.


Тогда $S_d(\ov V)=S_d(\ov w)$;


условие максимума (\ref{eq:(5.9)}) примет вид


$$


{H}_{v_j}(t)+\psi_V(t)


\begin{cases}


\leq 0 & \forall \, t\in [\ov\theta_0,\ov\theta_n]; \\


=0, & t\in S_{ac}(\ov{ w_j}),


\end{cases}


\qquad j=\ov{1,d(v)},


$$


где $S_{ac}(\ov{ w})$ --- множество


моментов времени, на котором сосредоточена абсолютно непрерывная


составляющая меры $\ov{ w}$; условие максимума (\ref{eq:(5.10)})


преобразуется к следующему:


$$


{H}_{v_j} \big|_{


\Theta(\ov{s}_i)}(\tau)+p_V^i(\tau)


\begin{cases}


\leq 0 & \forall \, \tau\in [0,\ov{d}_i]; \\


=0, & \ov{\omega}^i_j(\tau)>0,


\end{cases}\qquad j=\ov{1,d(v)}.


$$


При этом, если $p_V^i(\tau)<0$, то


для таких $\tau$ $\ov{\omega}^i(\tau)=1$; при строгом


неравенстве на $[0,\ov d_i]$ это равносильно равенству


$\ov{d}_i=|\ov{c^i}|$.





б) В случае $K=R^{d(v)}$ условия максимума также расписываются


покомпонентно: условие (\ref{eq:(5.9)}) примет вид


$$


\big|{H}_{v_j}(t)\big|+\psi_V(t)


\begin{cases}


\leq 0 & \forall \, t\in [\ov\theta_0,\ov\theta_n]; \\


=0, & t\in S_{ac}(\ov{ w_j}),


\end{cases}\qquad j=\ov{1,d(v)};


$$


в моменты $t\in S_{ac}(\ov{ w_j})$ знак ${H}_{v_j}(t)$ совпадает


со знаком абсолютно непрерывной составляющей $\ov


v_j(t)$; условие (\ref{eq:(5.10)}) преобразуется к следующему:


$$


\big|{H}_{v_j}|_{ \Theta(\ov{s}_i)}(\tau)\big|+p_V^i(\tau)


\begin{cases}


\leq 0 & \forall \; \tau\in [0,\ov{d}_i]; \\


=0, & \ov{\omega}^i_j(\tau)\ne0,


\end{cases} \qquad j=\ov{1,d(v)};


$$


знак ${H}_{v_j}\big|_{ \Theta(\ov{s}_i)}(\tau)$


совпадает со знаком предельного управления


$\ov{\omega}^i_j(\tau)$ на множестве, где последнее отлично от


нуля.





\section{Иллюстрирующие примеры}





Следующий модифицированный пример из


\cite{204} не только иллюстрирует ПМ, но и показывает, что при


аппроксимации обобщенной траектории обычными траекториями промежуточные


фазоограничения и их аналоги в предельных подсистемах могут не


выдерживаться.





\begin{exam}


Пусть управляемая система описывается уравнениями


\begin{equation}


\begin{array}{ll} \Dot x_1(t)=x_2(t)v(t), & x_1(0)=0, \\ [2mm]


\Dot x_2(t)=-x_1(t)v(t), & x_2(0)=1.


\end{array}\label{eq:(5.47)}


\end{equation}


Ставится задача минимизации функционала


$J=(x_2(\theta)+1)^2$ на множестве об. решений системы


(\ref{eq:(5.47)}) при ограничениях


\begin{gather}


v(t)\geq 0,\quad \int_0^1v(t)dt\leq \pi,\notag \\


\theta\in [0,1],\quad x_1(\theta-)\leq 0,\quad x_1(\theta


)\leq 0. \label{eq:(5.48)}


\end{gather}


(В \cite{204} рассмотрен этот пример с фазовым ограничением


$x_1(t)\leq 0$.)


Если не учитывать ограничение $z_1(\tau)\leq 0$ в предельной системе


(как в постановке задачи ${\cal P}(\theta)$), то оптимальным будет


решение


$$


\ov V(t)=


\begin{cases}


0, & t<\ov\theta; \\ \pi, & t\geq \ov\theta,


\end{cases}


\quad


\ov x_1(t)\equiv 0,\quad \ov x_2(t)=


\begin{cases}


1, & t<\ov\theta; \\ -1, & t\geq \ov\theta,


\end{cases}


\quad d\ov w=\pi\delta(t-\theta),


$$


где момент $\ov\theta$ --- любая точка из $(0,1)$, и $\min J=0$.





Рассмотрим предельную систему, соответствующую моменту


$\ov\theta$,


\begin{alignat*}{3}


\Dot z_1(\tau)&=z_2(\tau)\ov \omega(\tau),&\quad z_1(0)&=0,&\quad


z_1([\ov V(\ov\theta)])&=0, \\


%


\Dot z_2(\tau)&=-z_1(\tau)\ov\omega(\tau),&\quad z_2(0)&=1,&\quad


z_2([\ov V(\ov\theta)])&=-1.


\end{alignat*}


Поскольку $\ov \omega(\tau)\equiv 1$, то


$z_1(\tau)=\sin \tau$, $z_2(\tau)=\cos \tau.$


Очевидно, при $\tau \in (0,{\pi}/{2})$ \; $z_1(\tau)>0$.


Таким образом, хотя ограничение (\ref{eq:(5.48)}) выполняется на


оптимальном решении, соответствующее предельное решение,


определяющее скачок из $x(\ov\theta -)$ в


$x(\ov\theta)$, не удовлетворяет неравенству $z_1(\tau)\leq 0$.


Если оптимальную разрывную функцию $\ov


V(t)=\pi\chi_{_{[\theta,1]}}(t)$ аппроксимировать


последовательностью $\{V_n(t)\}$ такой, что $\pi/2\leq V_n(\theta)\leq


\pi$, то соответствующая последовательность обычных траекторий


$\{x_n\}$ будет удовлетворять ограничению (\ref{eq:(5.48)}); при


$0<V_n(\theta)<\pi/2$ ограничение (\ref{eq:(5.48)}) будет нарушаться.





Если же ввести в предельную систему фазовое ограничение $z_1(\tau)\leq


0$ $\forall\, \tau$, наследуемое от (\ref{eq:(5.48)}), то единственной


допустимой (и поэтому оптимальной) окажется траектория


$\wt x_1\equiv 0$, $\wt x_2\equiv 1$, $\wt V\equiv 0$, для


которой $\wt J=4 >\min J$.


\end{exam}








\begin{exam}


Требуется минимизировать функционал $J(d w)=x_3(20)-\nu x_2(20)$


на множестве обобщенных процессов


$(x(\cdot),d w)$ нелинейной управляемой системы


\begin{gather*}


\Dot


x_1=|x_2|,\quad \Dot x_2=1-v,\quad \Dot


x_3=x_1, \\


x_1(0)=1,\quad x_2(0)=5,\quad x_3(0)=0, \\


\theta = 10,\quad x_1(\theta)\geq 19,\quad v(t)\geq 0.


\end{gather*}





В зависимости от значения параметра $\nu$


ПМ выделяет следующие экстремали:





1) $\nu\geq 200\Rightarrow d\ov w\equiv 0$;








2) $\nu\in [128,200) \Rightarrow d\ov


w(t)=-5\delta(t)+\chi_{_{[0,\tau]}}dt,$ где $\tau=20-\sqrt{2\nu}$;





3) $\nu\in [50,128) \Rightarrow


d\ov w(t)=-5\delta(t)+\chi_{_{[0,4]}}dt;$





4) $\nu\in (0,50) \Rightarrow d\ov


w(t)=-5\delta(t)-6\delta(t-10)+\chi_{_{[0,4]}}dt+\chi_{_{[10,\tau]}}dt,$


где $\tau=20-\sqrt{2\nu}$.


\end{exam}





\section{Доказательство принципа максимума}





Доказательство теоремы 1


отличается от доказательства ПМ работы \cite{DS99} в деталях


расшифровки условий трансверсальности и оптимальности моментов, а


также --- в переходе к составной задаче. Последнее отличие состоит в том,


что начальной и конечной точкам $\theta_0$, $\theta_n$ сопоставляется


импульс (возможно, фиктивный). Это связано с наличием ограничений на


односторонние пределы $x(\theta_0)$, $V(\theta_0)$,


$x(\theta_n-)$, $V(\theta_n-)$.





Зафиксируем упорядоченное множество


моментов времени $\ov S$, включив в него точки разрыва функции $\ov


V$, промежуточные и граничные моменты, т.\,е.


$$


\ov{S}=\{\ov{s}_0,\dotsc,\ov{s}_N\} \supseteq


S_d(\ov V)\cup \{\ov \theta_0,\dotsc,\ov\theta_n\},


$$


где $\ov\theta_0=\ov{s}_0 \leq \ov{s}_1\leq\dotsb \leq


\ov{s}_N= \ov\theta_n.$ При этом, если вектор


$\ov \theta$ имеет совпадающие компоненты, то им соответствуют


столько же совпадающих точек множества $\ov S$, и других


совпадающих точек в $\ov S$ нет.





Пусть $\ov\omega^i\in


L_\infty([0,\ov d_i],K_1)$ произвольно,


$\ov d_i=[\ov V(s_i)]=0$ $\forall\, \ov


s_i\notin S_d(\ov V)$. Таким образом, все точки из


$\ov S$ делятся на три типа: 1)~$\ov s_i\in


S_d(\ov V)$, но $\ov s_i\neq \ov\theta_j$ $\forall \,


j\in \{0,\dotsc,n\}$; 2)~$\ov s_i\in S_d(\ov V)$ и


$\ov s_i=\ov\theta_j$ при некотором $j\in \{0,\dotsc,n\}$


(см. предположение (B7)); 3)~$\ov s_i\notin S_d(\ov V)$,


но может совпадать с некоторыми точками этого же типа, если им


соответствуют совпадающие точки из $\{\ov


\theta_0,\dotsc,\ov\theta_n\}$.





Обозначим через


$\ov x^i(\cdot)$, $\ov V^i(\cdot)$,


$\ov u^i(\cdot)$, $\ov v^i(\cdot)$


сужения функций $\ov x$, $\ov V$, $\ov u$, $\ov v$ на отрезки


$[\ov s_{i-1},\ov s_i]$, $i=\ov{1,N}$ (в правых концах


отрезков траекторные компоненты доопределяются по непрерывности, так


что $\ov x^i$, $\ov V^i\in AC$).





Перейдем к вспомогательной составной задаче ${\cal P}_N(\theta)$


аналогично доказательству ПМ работы \cite{DS99}. В задаче ${\cal


P}_N(\theta)$ имеются следующие ограничения:





уравнения динамики по $t$


\begin{gather*}


\Dot{x}^i=f(t,x^i ,V^i,u^i)+G(t,x^i,V^i)v^i,


\quad \Dot V^i=\|v^i\|,\\


u^i(t)\in U(t),\ \; v^i(t)\in K,


\quad t\in [s_{i-1},s_i], \quad i=\ov{1,N};


%\label{eq:(5.22)}


\end{gather*}





уравнения динамики по $\tau$


\begin{gather*}


\frac{dz^i}{d\tau}=G(s_i,z^i,z_V^i)\omega^i,\quad


\frac{dz_V^i}{d\tau}=1, \\


\omega^i(\tau)\in K_1,\quad\tau\in [0,d_i],\quad i=\ov{0,N}


% \label{eq:(5.23)}\end{equation}


\end{gather*}


(следовательно, в каждый


момент $s_i$ имитируется импульс), и дополнительные уравнения,


позволяющие избавиться от параметров в предельных подсистемах,


$$


\frac{ds_i}{d\tau}=0,\quad i=\ov{0,N};


$$





уравнения, связывающие концевые значения фазовых траекторий этапов,


\begin{gather}


x^{i+1}(s_i)-z^i(d_i)=0, \quad V^{i+1}(s_i)-z_V^i(d_i)=0,\notag \\


x^{i+1}(s_{i+1})-z^{i+1}(0)=0, \quad V^{i+1}(s_{i+1})-z_V^{i+1}(0)=0,


\quad i=\ov{0,N-1}; \label{eq:(5.63)} \\


s_0\leq s_1\leq \dotsb\leq s_N, \quad d_i\geq 0,\ \;


i=\ov{0,N}; \label{eq:(5.64)}


\end{gather}





ограничение, связывающее граничные значения моментов времени и


концы траекторий этапов,


\begin{equation}


\rho\in


C,\label{eq:(5.65)}


\end{equation}


где $\rho$ --- концевой вектор, который будет определен ниже.





Требуется минимизировать функционал $l(\rho)$.





Переход концевого вектора $b$ в


концевой вектор $\rho$ осуществляется по следующей


схеме (с переобозначением $t_0^i=s_{i-1}$, $t_1^i=s_i$ для


$i=\ov{1,N}$):


$$


\theta \rightarrow (t_0^1, (t_1^{i_k})_{i_k\in I} , t_1^{N}),


$$


\begin{alignat*}{2}


x(\theta-) &\rightarrow ( z^0(0),(


x^{i_k}(t_1^{i_k}))_{i_k\in I}, x^N(t_1^N) ), &\quad


x(\theta) &\rightarrow (x^1(t_0^1),


(x^{i_k+1}(t_0^{i_k+1}))_{i_k\in I} , z^{N}(d_N)), \\


%


V(\theta-) &\rightarrow ( z_V^0(0),(


V^{i_k}(t_1^{i_k}))_{i_k\in I}, V^N(t_1^N) ), &\quad


V(\theta) &\rightarrow (V^1(t_0^1),


(V^{i_k+1}(t_0^{i_k+1}))_{i_k\in I} , z_V^{N}(d_N)),


\end{alignat*}


где $I$ --- множество индексов (свое для каждого


допустимого процесса задачи ${\cal P}_N(\theta)$)


$$


I=\big\{1\leq i_k


\leq N-1,\ k=\ov{1,n-1} \mid i_k<i_{k+1} \big\}.


$$


Множество $I$ в совокупности с номерами $i=0,N$ определяет $n+1$


номеров этапов, граничные моменты времени которых и значения концов


траекторий входят в функционал и ограничение (\ref{eq:(5.65)}).





Ясно, что процесс $\ov e$ естественным образом


порождает мультипроцесс $\ov \xi$, допустимый в задаче ${\cal


P}_N(\theta)$ и являющийся в ней оптимальным. При этом


вектор $\ov b$ переходит в концевой


вектор $\ov\rho$ (условная запись $\ov b\rightarrow \ov\rho$)


по следующей схеме:


$$\ov\theta \rightarrow (\ov s_0,(\ov s_i)_{i\in \ov I},\ov s_N),


$$


\begin{alignat*}{2}


\ov x(\ov\theta-) &\rightarrow (\ov


z^0( 0),(\ov x^{i}(\ov s_i))_{i\in \ov I},\ov


x^N(\ov s_N)), &\quad \ov x(\ov\theta)


&\rightarrow (\ov x^1(\ov s_0),(\ov


x^{i+1}(\ov s_i))_{i\in \ov I},\ov z^{N}(\ov


d_N)), \\


%


\ov V(\ov\theta-) &\rightarrow


(\ov z_V^0( 0),(\ov V^{i}(\ov s_i))_{i\in \ov


I},\ov V^N(\ov s_N)), &\quad \ov


V(\ov\theta) &\rightarrow (\ov V^1(\ov


s_0),(\ov V^{i+1}(\ov s_i))_{i\in \ov I},\ov


z_V^N(\ov d_N)),


\end{alignat*}


где $$


\ov


I=\big\{1\leq i_k\leq N-1,\ i_k<i_{k+1}\mid \exists \,


\ov\theta_j = \ov s_{i_k} \big\}.


$$


(Множество $\ov I$ в объединении с $\{0,N\}$ --- это множество


номеров точек из $\ov S$, соответствующих моментам $\ov


\theta_j$,\ $j=\ov{0,n}$.)





Заменим ограничение $v^i(t)\in K$ на


$v^i(t)\in K_j =K\cap B_j$, где $j>\|\ov v(\cdot)\|_\infty$.


В результате получим семейство (по $j$) многоэтапных


задач ${\cal P}_{Nj}(\theta)$, в допустимые множества которых


исследуемый мультипроцесс $\ov \xi$ вкладывается очевидным образом и


является в этих задачах оптимальным.


К каждой задаче ${\cal P}_{Nj}(\theta)$ применим ПМ работы


\cite{Clarke.Vinter} и для доказательства теоремы 1 остается


расшифровать его условия.


Специфичной по сравнению


с доказательством ПМ из \cite{DS99} будет только расшифровка условия


трансверсальности ПМ для задачи ${\cal P}_{Nj}(\theta)$, на которой


мы и остановимся в данной работе.





В иной форме ограничения (\ref{eq:(5.63)})--(\ref{eq:(5.65)})


на граничные значения моментов времени и концы


траекторий этапов в задаче ${\cal P}_{Nj}(\theta)$ можно


записать в виде


\begin{gather*}


\Big(\rho,\big(z^{i_k}(0),z^{i_k}(d_{i_k}),z_V^{i_k}(0),


z_V^{i_k}(d_{i_k}), s_{i_k}(0),t_0^{i_k+1}\big)_{i_k\in I},


(z^0(d_0),z_V^0(d_0),s_0(0)),


(z^N(0),z_V^N(0),s_N(0))\Big)\in Q; \\


%


x^{i+1}(t_0^{i+1})-z^i(d_i)=0, \quad x^{i}(t_1^{i})-z^i(0)=0, \\


V^{i+1}(t_0^{i+1})-z_V^i(d_i)=0, \quad


V^{i}(t_1^{i})-z_V^i(0)=0,\\


%


(s_i(0),t_0^{i+1},t_1^i)\in


\{(a,a,a)\mid a\in R\}, \quad i\in \{1,\dotsc,N-1\}\setminus I; \\


%


\tau_0^i=0,\quad d_i\geq 0,\quad i=\ov{0,N}.


\end{gather*}


Здесь множество $Q$ определяется условиями


\begin{gather*}


\rho\in C;\\


x^{i_k+1}(t_0^{i_k+1})-z^{i_k}(d_{i_k})=0, \quad


V^{i_k+1}(t_0^{i_k+1})-z_V^{i_k}(d_{i_k})=0, \quad i_k\in


I\cup\{0\};\\


x^{i_k}(t_1^{i_k})-z^{i_k}(0)=0, \quad


V^{i_k}(t_1^{i_k})-z_V^{i_k}(0)=0,\quad i_k\in I\cup\{N\};\\


t_1^{i_k}=t_0^{i_k+1}=s_{i_k}(0),\quad i_k\in I, \quad t_0^1=


s_0(0),\quad s_N(0) =t_1^{N}.


\end{gather*}





Расшифровка условий трансверсальности для этапов, номера которых не


входят в множество $\ov I\cup\{0,N\}$ (и, следовательно, граничные


значения моментов времени и концы траекторий, соответствующие этим


этапам, не входят в функционал $l$ и не связаны множеством $C$,


задающим многоточечные фазоограничения), проводится так же, как в


доказательстве ПМ в работе \cite{DS99}. Если $i\notin \ov I\cup


\{0,N\}$, то


\begin{gather*}


\psi^{i+1}(\ov s_i)-p^i(\ov d_i)=0,\quad -\psi^i(\ov


s_i)+p^i(0)=0,\\


\psi_V^{i+1}(\ov s_i)-p_V^i(\ov d_i)=0,\quad -\psi_V^i(\ov


s_i)+p_V^i(0)=0,\\


p_s^i(0)-h_0^{i+1}+h_1^i=0,


\end{gather*}


причем $p_s^i(\ov d_i)=0$, $i=\ov{0,N}.$


При этом $\forall\, i=\ov{1,N}$ числа


$h_0^i$, $h_1^i$ удовлетворяют включениям


$$


h_0^{i+1}\in \co \ess_{t\rightarrow


\ov s_i}{\cal H}_0(t,\ov s_i+),\quad h_1^{i}\in


\co\ess_{t\rightarrow \ov


s_i}{\cal H}_0(t,\ov s_i-).


$$





Условия трансверсальности относительно вектора $\ov


\rho$ с учетом структуры множества $Q$ имеют вид


\begin{multline}


\label{eq:(5.66)}


\bigg(\Big((-h_0^1+p_s^0(0)),\


(p_s^i(0)+h_1^i-h_0^{i+1})_{i\in


\ov I}, \ (h_1^{N}+p_s^N(0)) \Big), \\


%


\Big( p^0(0), \ (-\psi^i(\ov s_i)+ p^i(0))_{i\in


\ov I}, \ (-\psi^N(\ov s_N)+ p^N(0)) \Big), \\


%


\Big((\psi^1(\ov s_0)- p^0(\ov d_0)), \ (\psi^{i+1}(\ov s_i)-


p^i(\ov d_i))_{i\in \ov I}, \ -p^N(\ov d_N)\Big), \\


%


\Big( p_V^0(0), \ (-\psi_V^i(\ov s_i)+


p_V^i(0))_{i\in \ov I}, \ (-\psi_V^N(\ov s_N)+ p_V^N(0))\Big),\\


%


\Big((\psi_V^1(\ov s_0)-


p_V^0(\ov d_0)), \ (\psi_V^{i+1}(\ov s_i)-


p_V^i(\ov d_i))_{i\in \ov I}, \ -p_V^N(\ov d_N)


\Big) \bigg)


\in \alpha_0\partial l(\ov \rho)+N_C(\ov \rho).


\end{multline}





Введем обозначения


$$


h_{s_0}^{+}=h_0^1,\quad


h_{s_N}^{-}=h_1^N,\quad h_{s_i}^{+}=h_0^{i+1},\quad


h_{s_i}^{-}=h_1^{i},\quad i=\ov{1,N-1},


$$


и положим $\psi(\ov s_{i-1}+)=\psi^{i}(t_0^i)$, $\psi(\ov s_i


-)=\psi^{i}(t_1^{i})$, $\psi_V(\ov s_{i-1}+)=\psi_V^{i}(t_0^i)$,


$\psi_V(\ov s_i -)=\psi_V^{i}(t_1^{i})$, $i=\ov{1,N}$.


Тогда из (\ref{eq:(5.66)}) с учетом интегрального


представления $p_s^i(0)=\int\limits_0^{\ov


d_i}H_{1t}\big|_{ \Theta(\ov s_i)}(\tau)d\tau$ и переобозначения


$\ov b\rightarrow \ov\rho$ получаем условия


трансверсальности и оптимальности моментов теоремы 1.$\quad\square$


\medskip





Докажем замечания, следующие за теоремой 1.





1) Приведем схему перехода к вспомогательной


многоэтапной задаче и расшифровку условий трансверсальности, если в


задаче ${\cal P}(\theta)$ отсутствуют ограничения на значения


$x(\theta_0)$, $V(\theta_0)$, $x(\theta_n-)$, $V(\theta_n-)$ и эти


значения не входят в функционал $l$. Теперь начальному и конечному


моментам времени не сопоставляется обязательно импульс, как и при


доказательстве ПМ в \cite{DS99} ($\ov\theta_0\leq \ov


s_0$, $\ov s_N\leq \ov \theta_n$),


и вектор $\ov b$ переходит в концевой вектор $\ov\rho$ следующим


образом:


$$


\ov\theta \rightarrow (\ov \theta_0,(\ov


s_i)_{i\in \ov I},\ov \theta_n),


$$


\begin{alignat*}{2}


\ov x(\ov\theta-)


&\rightarrow (\ov x^0(\ov \theta_0),(\ov x^{i}(\ov s_i))_{i\in


\ov I}), &\quad \ov x(\ov\theta) &\rightarrow ((\ov x^{i+1}(\ov


s_i))_{i\in \ov I},\ov x^{N+1}(\ov \theta_n)), \\


\ov V(\ov\theta-) &\rightarrow (\ov V^0(\ov \theta_0),(\ov V^{i}(\ov


s_i))_{i\in \ov I}), &\quad \ov V(\ov\theta) &\rightarrow ((\ov


V^{i+1}(\ov s_i))_{i\in \ov I},\ov V^{N+1}(\ov \theta_n)).


\end{alignat*}


Для допустимого процесса задачи ${\cal P}_{N}(\theta)$


переход вектора $b$ в концевой вектор $\rho$ происходит по следующей


схеме (с переобозначением $t_0^i=s_{i-1}$, $t_1^i=s_i$ для


$i=\ov{0,N+1}$):


$$


\theta \rightarrow ( t_0^0,


(t_1^{i_k})_{i_k\in I} , t_1^{N+1} ),


$$


\begin{alignat*}{2}


x(\theta-) &\rightarrow ( x^0(t_0^0),(


x^{i_k}(t_1^{i_k}))_{i_k\in I} ), &\quad x(\theta)


&\rightarrow ( (x^{i_k+1}(t_0^{i_k+1}))_{i_k\in I},


x^{N+1}(t_1^{N+1})), \\


V(\theta-) &\rightarrow (


V^0(t_0^0),( V^{i_k}(t_1^{i_k}))_{i_k\in I} ), &\quad


V(\theta) &\rightarrow (


(V^{i_k+1}(t_0^{i_k+1}))_{i_k\in I}, V^{N+1}(t_1^{N+1})).


\end{alignat*}





Ограничения на граничные значения моментов времени и концы траекторий


этапов в задаче ${\cal P}_{N}(\theta)$ задаются условиями


\begin{gather*}


\Big(\rho,\big(z^{i_k}(0),z^{i_k}(d_{i_k}),z_V^{i_k}(0),


z_V^{i_k}(d_{i_k}) ,


s_{i_k}(0),t_0^{i_k+1}\big)_{i_k\in I},s_0(0),s_N(d_N)\Big)\in Q;\\


%


x^{i+1}(t_0^{i+1})-z^i(d_i)=0,\quad


x^{i}(t_1^{i})-z^i(0)=0, \\


%


V^{i+1}(t_0^{i+1})-z_V^i(d_i)=0,\quad


V^{i}(t_1^{i})-z_V^i(0)=0,\quad  i\in


\{0,\dotsc,N\}\setminus I; \\


%


(s_i(0),t_0^{i+1},t_1^i)\in \{(a,a,a)\mid a\in R\}, \quad


i\in \{1,\dotsc,N-1\}\setminus I; \\


%


\tau_0^i=0,\quad d_i\geq


0,\quad i=\ov{0,N} .


\end{gather*}


Здесь множество $Q$


определяется условиями 


\begin{gather*}


\rho\in C;\quad x^{i_k+1}(t_0^{i_k+1})=z^{i_k}(d_{i_k}),\quad


V^{i_k+1}(t_0^{i_k+1})=z_V^{i_k}(d_{i_k}),\\


x^{i_k}(t_1^{i_k})=z^{i_k}(0),\quad


V^{i_k}(t_1^{i_k})=z_V^{i_k}(0),\quad i_k\in I; \\


t_1^{i_k}=t_0^{i_k+1}=s_{i_k}(0),\quad i_k\in I\cup\{0,N\},\\


t_0^0\leq s_0(0),\quad s_N(0)\leq t_1^{N+1}.


\end{gather*}





Условия трансверсальности относительно вектора $\ov


\rho$ с учетом структуры множества $Q$ принимают вид


\begin{multline*}


\bigg(\Big((-h_0^0-a_0),\ 


(p_s^i(0)+h_1^i-h_0^{i+1})_{i\in \ov I}, \ (h_1^{N+1}+a_1) \Big),\


\Big( \psi^0(\ov\theta_0), \ (-\psi^i(\ov s_i)+ p^i(0))_{i\in


\ov I} \Big), \\


%


\Big((\psi^{i+1}(\ov s_i)- p^i(\ov


d_i))_{i\in \ov I},\ -\psi^{N+1}(\ov \theta_n) \Big), \ \Big(


\psi_V^0(\ov\theta_0), (-\psi_V^i(\ov s_i)+ p_V^i(0))_{i\in \ov I}


\Big), \\


%


\Big( (\psi_V^{i+1}(\ov s_i)- p_V^i(\ov


d_i))_{i\in \ov I}, -\psi_V^{N+1}(\ov \theta_n) \Big) \bigg) \in


\alpha_0\partial l(\ov \rho)+N_C(\ov \rho),


\end{multline*}


где $a_0$, $a_1\geq 0$. При этом


$$


p_s^0(0)+h_1^0-h_0^1=-a_0, \quad p_s^N(0)+h_1^{N}-h_0^{N+1}=a_1.


$$


Так как $\ov\theta_0=\ov s_0$,


$\ov\theta_n=\ov s_N$, то согласно


замечанию 2) к теореме 2 из \cite{DS99} $h_0^0=h_1^0$,


$h_0^{N+1}=h_1^{N+1}.$ Следовательно,


$$


-h_0^0-a_0=p_s^0(0)-h_0^1,\quad h_1^{N+1}+a_1=p_s^N(0)+h_1^N.


$$


Тем самым доказано замечание 1), следующее после теоремы 1.





2) Рассмотрим случай совпадения промежуточных моментов.


По предположению (B7), если промежуточные моменты совпадают, то


эта точка не является точкой импульса. Пусть


$$


\ov\theta_j=\ov\theta_{j+1}=\dotsb=\ov\theta_{j+k},


$$


и этим точкам соответствуют совпадающие моменты


$\ov s_m, \ov s_{m+1}, \dotsc,$


$\ov s_{m+k}$. Тогда


$$


h_0^{i}=h_1^{i},\quad i=\ov{m+1,m+k}, \quad \ov


d_i=0, \quad p_s^i(0)=0,\quad p^i(0)=p^i(\ov d_i), \quad


i=\ov{m,m+k}.


$$








\begin{thebibliography}{9}


\bibitem{80}


Миллер Б.М. {\it Оптимизация динамических систем с


обобщенным управлением} // Автоматика и телемеханика. -- 1989. --


\No\,6. -- С.\,23--34.


\bibitem{75}


Завалищин С.Т., Сесекин А.Н.


{\it Импульсные процессы: модели и приложения.} -- М.:


Наука, 1991. -- 256~с.


\bibitem{83}


Vinter R.B., Pereira F.M.


{\it A maximum principle for optimal processes with


discontinuous trajectories} // SIAM J. Contr. and Optim. --


1988. -- V.\,26. -- \No\,1. -- P.\,205--229.


\bibitem{Diz.vuz}


Дыхта В.А.


{\it Необходимые условия оптимальности импульсных процессов при


ограничениях на образ управляющей меры} // Изв. вузов. Математика. --


1996. -- \No\,12. -- С.\,1--9.


\bibitem{DS99}


Дыхта В.А., Самсонюк О.Н.


{\it Принцип максимума в негладких задачах оптимального управления с


разрывными траекториями} // Изв. вузов. Математика. -- 1999. -- \No\,12.


-- С.\,26--37.


\bibitem{DIFAC} { Dykhta V. } {\it Optimality conditions for


impulsive processes and its applications} // Prepr. of 13th World


Congress IFAC, 1996. -- V.\,D. -- P.\,345--350.


\bibitem{Clarke.Vinter}


Clarke F.H., Vinter R.B. {\it


Optimal multiprocesses} // SIAM J. Contr. and Optim. -- 1989. -- V.\,27.


-- \No\,5. -- P.\,1072--1091.


\bibitem{23}


Кларк Ф. {\it Оптимизация и негладкий анализ.} -- М.: Наука, 1988.


-- 280~с.


\bibitem{204}


Миллер Б.М. {\it Условия


оптимальности в задачах обобщенного управления} // Автоматика и


телемеханика. -- 1992. -- \No\,5. -- С.\,50--58.





\end{thebibliography}





\vskip 2cm





\post{\it Иркутская государственная}{\it Поступила}


\post{\it экономическая академия}{31.07.2000}





\label{end}


\end{document}





